\chapter{Conclusion and Perspectives}
%%%%%%%%%%%%%%%%%%%%%%%%%%%%%%%%%%%%%%%%%%%%%%%%%%%%%%%%%%%%%%%%%%%%%%%%%%%%%%%%
\section{Conclusion}
Both validations using spectralon with albedo of $0.5$ and comparing measured data to the ones measured by UGE, prove that the proposed spectral BRDF measurement approached shows a good accuracy.
For the study of spectral BRDF of roads surface, we conduct many measurements on the multi-material roads and their components.
The data demonstrates that the binder has a significant contribution on the BRDF behavior of the mixture roads.
Besides, the aggregates with different size has few difference except at the specular direction, which may have different contribution on the BRDF of mixture at the specular direction.


%%%%%%%%%%%%%%%%%%%%%%%%%%%%%%%%%%%%%%%%%%%%%%%%%%%%%%%%%%%%%%%%%%%%%%%%%%%%%%%%
\section{Perspectives}

In future work, the measured spectral BRDF data can be used to calibrate and evaluate existing reflectance models, such as the Cook Torrance model~\cite{1982_Cook}, the Oren Nayar model~\cite{1995_Oren}, and the Ashikhmin Shirley model~\cite{2000_Ashikhmin}. 
Furthermore, it will be of interest to develop a hybrid model that combines the most appropriate BRDF models for each material component to represent the BRDF of composite road surfaces.
This can provide a foundation for predicting the spectral BRDF of more complex pavement materials.